% GENERATED FILE. Edit canonical lessons, then run npm run generate:books.
% canonical-source-hash: fnv1a64:7acafcc5ebd01b32
% canonical-lessons: TA-C341-akalvaraycci, TA-C341-tolporul, TA-C341-putaiyal, TA-C341-tontiyetu, TA-C341-kantetu

\chapter{Digging Up the Past}
\label{ch:ta-digging-up-the-past}
\hlchaptermodality{341}

\begin{chapteropening}
\textbf{By the end of this chapter:} \emph{I can talk about digs and the things found in them.}

Complete the last lesson of chapter 341: I can talk about digs and the things found in them.
\end{chapteropening}

\section[\texorpdfstring{\ta{அகழ்வாராய்ச்சி} (akaḻvārāycci)}{அகழ்வாராய்ச்சி (akaḻvārāycci)}]{\ta{அகழ்வாராய்ச்சி} (akaḻvārāycci) — an excavation}

\label{lesson:TA-C341-akalvaraycci}



\begin{quote}
Before the new one: say the Tamil for a golden age, then the Tamil for a civilization.
\end{quote}

\subsection*{You'll want to know: \ta{அகழ்வாராய்ச்சி}}
\textbf{\ta{அகழ்வாராய்ச்சி}} — \emph{akaḻvārāycci} — \textquotedblleft{}an excavation\textquotedblright{}.

\begin{quote}
\ta{அகழ்வாராய்ச்சியில்} \ta{பழைய} \ta{நாணயங்கள்} \ta{கிடைத்தன}.
\end{quote}

\begin{quote}
\emph{akaḻvārāycciyil paḻaiya nāṇayaṅkaḷ kiṭaittaṉa.}
\end{quote}

\begin{quote}
Old coins were found in the excavation.
\end{quote}

A dig made to learn about the past; \ta{அகழ்}, to dig, is in it.

\begin{cousinweb}
\textbf{\ta{அகழ்வாராய்ச்சி}} (\emph{akaḻvārāycci}) is a Tamil compound built from \ta{அகழ்வு}, \textquotedblleft{}digging\textquotedblright{}, from \ta{அகழ்} (DEDR 11), and \ta{ஆராய்ச்சி}, \textquotedblleft{}research\textquotedblright{} (DEDR 377).
\end{cousinweb}

\begin{grammarlens}[title={What you've built}]
One more word for telling what happened.
\end{grammarlens}

\subsection*{Your turn}
\begin{itemize}
\raggedright
  \item \emph{Say it:} \emph{akaḻvārāycci}
  \item \emph{Say it:} \emph{akaḻvārāycci}, once more
  \item \emph{Say it:} \emph{akaḻvārāycciyil paḻaiya nāṇayaṅkaḷ kiṭaittaṉa.} — Old coins were found in the excavation.
  \item \emph{Recall:} say the Tamil for a golden age, then the Tamil for a civilization, then say \emph{akaḻvārāycci} again
\end{itemize}

\begin{tcolorbox}[breakable,colback=teal!4,colframe=teal!35!black,title={Before you move on}]
What does \ta{அகழ்வாராய்ச்சி} mean? (\textquotedblleft{}An excavation\textquotedblright{}.) Say it once more.
\end{tcolorbox}
\section[\texorpdfstring{\ta{தொல்பொருள்} (tolporuḷ)}{தொல்பொருள் (tolporuḷ)}]{\ta{தொல்பொருள்} (tolporuḷ) — an antiquity; an ancient object}

\label{lesson:TA-C341-tolporul}



\begin{quote}
Before the new one: say the Tamil for a civilization, then the Tamil for an excavation.
\end{quote}

\subsection*{You'll want to know: \ta{தொல்பொருள்}}
\textbf{\ta{தொல்பொருள்}} — \emph{tolporuḷ} — \textquotedblleft{}an antiquity; an ancient object\textquotedblright{}.

\begin{quote}
\ta{அருங்காட்சியகத்தில்} \ta{பல} \ta{தொல்பொருள்கள்} \ta{இருக்கின்றன}.
\end{quote}

\begin{quote}
\emph{aruṅkāṭciyakattil pala tolporuḷkaḷ irukkiṉṟaṉa.}
\end{quote}

\begin{quote}
There are many antiquities in the museum.
\end{quote}

\ta{தொல்}, ancient + \ta{பொருள்}, thing.

\begin{cousinweb}
\textbf{\ta{தொல்பொருள்}} (\emph{tolporuḷ}) is a Tamil compound built from \ta{தொல்}, \textquotedblleft{}old, ancient\textquotedblright{} (DEDR 3516), and \ta{பொருள்}, \textquotedblleft{}thing\textquotedblright{} (DEDR 4544).
\end{cousinweb}

\begin{grammarlens}[title={What you've built}]
One more word for telling what happened.
\end{grammarlens}

\subsection*{Your turn}
\begin{itemize}
\raggedright
  \item \emph{Say it:} \emph{tolporuḷ}
  \item \emph{Say it:} \emph{tolporuḷ}, once more
  \item \emph{Say it:} \emph{aruṅkāṭciyakattil pala tolporuḷkaḷ irukkiṉṟaṉa.} — There are many antiquities in the museum.
  \item \emph{Recall:} say the Tamil for a civilization, then the Tamil for an excavation, then say \emph{tolporuḷ} again
\end{itemize}

\begin{tcolorbox}[breakable,colback=teal!4,colframe=teal!35!black,title={Before you move on}]
What does \ta{தொல்பொருள்} mean? (\textquotedblleft{}An antiquity; an ancient object\textquotedblright{}.) Say it once more.
\end{tcolorbox}
\section[\texorpdfstring{\ta{புதையல்} (putaiyal)}{புதையல் (putaiyal)}]{\ta{புதையல்} (putaiyal) — a treasure; buried treasure}

\label{lesson:TA-C341-putaiyal}



\begin{quote}
Before the new one: say the Tamil for an excavation, then the Tamil for an antiquity.
\end{quote}

\subsection*{You'll want to know: \ta{புதையல்}}
\textbf{\ta{புதையல்}} — \emph{putaiyal} — \textquotedblleft{}a treasure; buried treasure\textquotedblright{}.

\begin{quote}
\ta{ஒரு} \ta{விவசாயி} \ta{தன்} \ta{நிலத்தில்} \ta{புதையலைக்} \ta{கண்டுபிடித்தார்}.
\end{quote}

\begin{quote}
\emph{oru vivacāyi taṉ nilattil putaiyalaik kaṇṭupiṭittār.}
\end{quote}

\begin{quote}
A farmer found treasure in his field.
\end{quote}

From \ta{புதை}, to bury: buried treasure.

\begin{cousinweb}
\textbf{\ta{புதையல்}} (\emph{putaiyal}) is native Dravidian: DEDR 4509 gives it under \ta{புதை}, \textquotedblleft{}to be buried\textquotedblright{}, as \textquotedblleft{}hidden treasure\textquotedblright{}.
\end{cousinweb}

\begin{grammarlens}[title={What you've built}]
One more word for telling what happened.
\end{grammarlens}

\subsection*{Your turn}
\begin{itemize}
\raggedright
  \item \emph{Say it:} \emph{putaiyal}
  \item \emph{Say it:} \emph{putaiyal}, once more
  \item \emph{Say it:} \emph{oru vivacāyi taṉ nilattil putaiyalaik kaṇṭupiṭittār.} — A farmer found treasure in his field.
  \item \emph{Recall:} say the Tamil for an excavation, then the Tamil for an antiquity, then say \emph{putaiyal} again
\end{itemize}

\begin{tcolorbox}[breakable,colback=teal!4,colframe=teal!35!black,title={Before you move on}]
What does \ta{புதையல்} mean? (\textquotedblleft{}A treasure; buried treasure\textquotedblright{}.) Say it once more.
\end{tcolorbox}
\section[\texorpdfstring{\ta{தோண்டியெடு} (tōṇṭiyeṭu)}{தோண்டியெடு (tōṇṭiyeṭu)}]{\ta{தோண்டியெடு} (tōṇṭiyeṭu) — to dig out; to unearth}

\label{lesson:TA-C341-tontiyetu}



\begin{quote}
Before the new one: say the Tamil for an antiquity, then the Tamil for a treasure.
\end{quote}

\subsection*{You'll want to know: \ta{தோண்டியெடு}}
\textbf{\ta{தோண்டியெடு}} — \emph{tōṇṭiyeṭu} — \textquotedblleft{}to dig out; to unearth\textquotedblright{}.

\begin{quote}
\ta{அவர்கள்} \ta{ஒரு} \ta{பழைய} \ta{பானையைத்} \ta{தோண்டியெடுத்தார்கள்}.
\end{quote}

\begin{quote}
\emph{avarkaḷ oru paḻaiya pāṉaiyait tōṇṭiyeṭuttārkaḷ.}
\end{quote}

\begin{quote}
They dug out an old pot.
\end{quote}

Past (I …): \ta{தோண்டியெடுத்தேன்}, \emph{tōṇṭiyeṭuttēṉ}.

\ta{தோண்டு}, to dig + \ta{எடு}, to take out.

\begin{cousinweb}
\textbf{\ta{தோண்டியெடு}} (\emph{tōṇṭiyeṭu}) is a Tamil compound built from \ta{தோண்டு}, \textquotedblleft{}to dig, excavate\textquotedblright{} (DEDR 3549), and \ta{எடு}, \textquotedblleft{}to take up\textquotedblright{} (DEDR 851).
\end{cousinweb}

\begin{grammarlens}[title={What you've built}]
One more word for telling what happened.
\end{grammarlens}

\subsection*{Your turn}
\begin{itemize}
\raggedright
  \item \emph{Say it:} \emph{tōṇṭiyeṭu}
  \item \emph{Say it:} \emph{tōṇṭiyeṭu}, once more
  \item \emph{Say it:} \emph{avarkaḷ oru paḻaiya pāṉaiyait tōṇṭiyeṭuttārkaḷ.} — They dug out an old pot.
  \item \emph{Recall:} say the Tamil for an antiquity, then the Tamil for a treasure, then say \emph{tōṇṭiyeṭu} again
\end{itemize}

\begin{tcolorbox}[breakable,colback=teal!4,colframe=teal!35!black,title={Before you move on}]
What does \ta{தோண்டியெடு} mean? (\textquotedblleft{}To dig out; to unearth\textquotedblright{}.) Say it once more.
\end{tcolorbox}
\section[\texorpdfstring{\ta{கண்டெடு} (kaṇṭeṭu)}{கண்டெடு (kaṇṭeṭu)}]{\ta{கண்டெடு} (kaṇṭeṭu) — to come upon; to find}

\label{lesson:TA-C341-kantetu}



\begin{quote}
Before the new one: say the Tamil for a treasure, then the Tamil for to dig out.
\end{quote}

\subsection*{You'll want to know: \ta{கண்டெடு}}
\textbf{\ta{கண்டெடு}} — \emph{kaṇṭeṭu} — \textquotedblleft{}to come upon; to find\textquotedblright{}.

\begin{quote}
\ta{சிறுவன்} \ta{ஆற்றங்கரையில்} \ta{ஒரு} \ta{பழைய} \ta{நாணயத்தைக்} \ta{கண்டெடுத்தான்}.
\end{quote}

\begin{quote}
\emph{ciṟuvaṉ āṟṟaṅkaraiyil oru paḻaiya nāṇayattaik kaṇṭeṭuttāṉ.}
\end{quote}

\begin{quote}
The boy came upon an old coin on the river bank.
\end{quote}

Past (I …): \ta{கண்டெடுத்தேன்}, \emph{kaṇṭeṭuttēṉ}.

To find something by chance; news says \ta{கண்டெடுக்கப்பட்டது}, was found.

\begin{cousinweb}
\textbf{\ta{கண்டெடு}} (\emph{kaṇṭeṭu}) is a Tamil compound built from \ta{கண்டு}, \textquotedblleft{}having seen\textquotedblright{}, from \ta{காண்} (DEDR 1443), and \ta{எடு}, \textquotedblleft{}to take up\textquotedblright{} (DEDR 851).
\end{cousinweb}

\begin{grammarlens}[title={What you've built}]
One more word for telling what happened.
\end{grammarlens}

\subsection*{Your turn}
\begin{itemize}
\raggedright
  \item \emph{Say it:} \emph{kaṇṭeṭu}
  \item \emph{Say it:} \emph{kaṇṭeṭu}, once more
  \item \emph{Say it:} \emph{ciṟuvaṉ āṟṟaṅkaraiyil oru paḻaiya nāṇayattaik kaṇṭeṭuttāṉ.} — The boy came upon an old coin on the river bank.
  \item \emph{Recall:} say the Tamil for a treasure, then the Tamil for to dig out, then say \emph{kaṇṭeṭu} again
\end{itemize}

\begin{tcolorbox}[breakable,colback=teal!4,colframe=teal!35!black,title={Before you move on}]
What does \ta{கண்டெடு} mean? (\textquotedblleft{}To come upon; to find\textquotedblright{}.) Say it once more.
\end{tcolorbox}
